% Publication Materials for RAG Comparison Research
% Generated: January 10, 2026
% FINAL VERSION WITH ACTUAL RESULTS

\documentclass{article}
\usepackage{booktabs}
\usepackage{amsmath}
\usepackage{graphicx}
\usepackage{hyperref}

\title{Comprehensive Evaluation of RAG Knowledge Bases for\\Smoking Cessation Counseling}
\author{Research Team}
\date{January 2026}

\begin{document}

\maketitle

\begin{abstract}
We present a rigorous comparison of Retrieval-Augmented Generation (RAG) approaches for smoking cessation counseling, using a novel independent test set validated for zero data leakage. Our analysis includes multiple metrics (ROUGE-L, BLEU, METEOR) across five independent trials with statistical significance testing.

\textbf{Key Finding:} Contrary to expectations, the baseline approach (no retrieval) outperformed all RAG variants (ROUGE-L: 0.231 vs 0.224 for best RAG). Among RAG approaches, AI-generated knowledge bases slightly exceeded human-curated ones (102.1\% relative performance), meeting the first research target. However, the improvement over raw sources (+7.8\%) fell short of the 40\% target.

\textbf{Methodological Contribution:} We discovered severe data leakage in the original test set (7/15 questions with $\geq$90\% similarity to training data), invalidating previous results. This highlights the critical importance of test set independence verification in NLP research.
\end{abstract}

\section{Methodology}

\subsection{Data Leakage Discovery}
During our audit, we discovered that the original test set had significant overlap with training data:

\begin{table}[h]
\centering
\caption{Original Test Set Data Leakage Analysis}
\begin{tabular}{lcc}
\toprule
Severity Level & Count & Percentage \\
\midrule
Severe ($\geq$90\% similarity) & 7 & 47\% \\
Moderate (70--90\% similarity) & 6 & 40\% \\
Clean ($<$70\% similarity) & 2 & 13\% \\
\bottomrule
\end{tabular}
\label{tab:leakage}
\end{table}

Four test questions were \textit{exact matches} (100\% similarity) to questions in the training data, fundamentally invalidating the evaluation.

\subsection{Independent Test Set Creation}
We created a new 15-question test set with:
\begin{itemize}
    \item Maximum 50\% similarity to any training question (fuzzy string matching)
    \item Coverage of all original topics (benefits, cravings, NRT, relapse, etc.)
    \item Grammatically distinct formulations
    \item Domain expert-validated reference answers
\end{itemize}

\subsection{Evaluation Protocol}
\begin{itemize}
    \item \textbf{Model:} GPT-4o-mini (temperature=0.3, max\_tokens=200)
    \item \textbf{Retrieval:} ChromaDB with cosine similarity, k=3 documents
    \item \textbf{Runs:} 5 independent trials
    \item \textbf{Statistical Analysis:} 95\% confidence intervals, paired t-tests, Cohen's d
\end{itemize}

\section{Results}

\subsection{Main Results}

\begin{table}[h]
\centering
\caption{RAG Approach Comparison (5 runs, 15 questions each)}
\begin{tabular}{lcccc}
\toprule
Approach & ROUGE-L & BLEU & METEOR \\
\midrule
Baseline (No RAG) & $0.231 \pm 0.008$ & $0.054 \pm 0.004$ & $0.310 \pm 0.010$ \\
Raw Sources RAG & $0.208 \pm 0.004$ & $0.039 \pm 0.003$ & $0.275 \pm 0.008$ \\
AI-Generated RAG & $0.224 \pm 0.005$ & $0.057 \pm 0.003$ & $0.290 \pm 0.008$ \\
Human-Curated RAG & $0.220 \pm 0.004$ & $0.045 \pm 0.003$ & $0.296 \pm 0.010$ \\
\bottomrule
\end{tabular}
\label{tab:main_results}
\end{table}

Values shown as mean $\pm$ 95\% CI.

\subsection{Statistical Comparisons (ROUGE-L)}

\begin{table}[h]
\centering
\caption{Pairwise Statistical Comparisons}
\begin{tabular}{lcccc}
\toprule
Comparison & Difference & Cohen's d & p-value & Significant \\
\midrule
AI vs Human & +2.1\% & 0.72 & 0.184 & No \\
AI vs Raw & +7.8\% & 3.19 & 0.002 & Yes** \\
Human vs Raw & +5.5\% & 1.89 & 0.014 & Yes* \\
\bottomrule
\end{tabular}
\label{tab:statistical}
\end{table}

* $p < 0.05$, ** $p < 0.01$

\subsection{Professor's Targets Assessment}

\begin{table}[h]
\centering
\caption{Research Target Evaluation}
\begin{tabular}{lccc}
\toprule
Target & Expected & Actual & Status \\
\midrule
AI achieves 84\% of Human & $\geq$84\% & 102.1\% & \textbf{MET} \\
AI 40\% better than Raw & $\geq$+40\% & +7.8\% & NOT MET \\
\bottomrule
\end{tabular}
\label{tab:targets}
\end{table}

\section{Discussion}

\subsection{Key Findings}

\textbf{1. Baseline outperforms RAG:} The most surprising result is that the no-retrieval baseline achieved the highest ROUGE-L score (0.231), suggesting that for well-documented domains like smoking cessation, the LLM's parametric knowledge may be sufficient.

\textbf{2. AI-generated matches human-curated:} AI-generated knowledge bases performed comparably to human-curated ones (102.1\% relative performance), with no statistically significant difference ($p = 0.184$).

\textbf{3. Processed data beats raw sources:} Both AI-generated and human-curated RAG significantly outperformed raw sources ($p < 0.05$), confirming the value of structured knowledge bases.

\textbf{4. Data leakage reversed conclusions:} The original test set showed Raw $>$ Human $>$ AI, which completely reversed with proper independence validation.

\subsection{Why This Matters}
The discovery of data leakage in the original test set demonstrates:
\begin{enumerate}
    \item The importance of rigorous test set validation in NLP research
    \item That high performance scores may be artifacts of data contamination
    \item The need for fuzzy string matching during test set creation
\end{enumerate}

\subsection{Implications for RAG Research}

Our findings suggest a domain-dependency hypothesis: RAG benefit may be \textit{inversely related} to domain representation in LLM training data. For well-documented topics:
\begin{itemize}
    \item Baseline LLM knowledge may be comprehensive
    \item Retrieved context may introduce noise rather than signal
    \item RAG overhead may not justify marginal improvements
\end{itemize}

Future work should test this hypothesis across domains with varying LLM knowledge saturation.

\subsection{Limitations}
\begin{itemize}
    \item Single domain (smoking cessation) - results may not generalize
    \item Single LLM (GPT-4o-mini) - other models may show different patterns
    \item 15 test questions - larger test sets would provide more statistical power
    \item Domain is well-covered in LLM training data - may reduce RAG benefit
    \item Metrics measure lexical overlap, not semantic accuracy or factual correctness
\end{itemize}

\section{Conclusion}

This study makes two primary contributions:

\textbf{Methodological:} We demonstrate the critical importance of test set independence verification. Data leakage completely reversed our conclusions, serving as a warning for NLP researchers.

\textbf{Empirical:} We find that RAG does not universally improve performance. For well-known domains, LLM parametric knowledge may suffice. When RAG is used, AI-generated knowledge bases perform comparably to human-curated ones at a fraction of the cost.

These findings shift the research question from ``Does RAG help?'' to ``When does RAG help?'' - a more nuanced and practically valuable inquiry.

\section*{Acknowledgments}
We thank the professor for guidance and the research team for data curation efforts.

\appendix
\section{Effect Size Reference}

Cohen's d interpretation:
\begin{itemize}
    \item $|d| < 0.2$: Negligible
    \item $0.2 \leq |d| < 0.5$: Small
    \item $0.5 \leq |d| < 0.8$: Medium
    \item $|d| \geq 0.8$: Large
\end{itemize}

\section{Complete Metrics by Approach}

\begin{table}[h]
\centering
\caption{Complete Metrics (Mean $\pm$ 95\% CI)}
\begin{tabular}{lccccc}
\toprule
Approach & ROUGE-1 & ROUGE-2 & ROUGE-L & BLEU & METEOR \\
\midrule
Baseline & $0.342 \pm 0.010$ & $0.097 \pm 0.012$ & $0.231 \pm 0.008$ & $0.054 \pm 0.004$ & $0.310 \pm 0.010$ \\
Raw & $0.316 \pm 0.005$ & $0.068 \pm 0.004$ & $0.208 \pm 0.004$ & $0.039 \pm 0.003$ & $0.275 \pm 0.008$ \\
AI & $0.325 \pm 0.004$ & $0.086 \pm 0.004$ & $0.224 \pm 0.005$ & $0.057 \pm 0.003$ & $0.290 \pm 0.008$ \\
Human & $0.334 \pm 0.007$ & $0.079 \pm 0.005$ & $0.220 \pm 0.004$ & $0.045 \pm 0.003$ & $0.296 \pm 0.010$ \\
\bottomrule
\end{tabular}
\label{tab:complete}
\end{table}

\end{document}
